%%% ============================================================
%%% The First-G Event and a Discrete Sinc^2 Identity
%%%
%%% B.R. Sanders, M. Gish (2026)
%%% Target venue: Integers -- Electronic Journal of Combinatorial
%%%               Number Theory
%%% MSC: 11A41, 11N05, 11A51, 11Y05, 42A16
%%%
%%% DOI for verification scripts: 10.5281/zenodo.18852047
%%%
%%% Source: Fork A restoration of held draft
%%%   first_g_sinc2_FINAL.tex (private corpus, _held_first_g)
%%% Surgical additions:
%%%   - lens-ownership preamble (§0)
%%%   - PROVEN/COMPUTED/RHYME/OPEN tier-discipline paragraph (§1)
%%%   - bibliography expansion (Erdos, Pomerance, Tenenbaum,
%%%     Iwaniec-Kowalski, Friedlander-Iwaniec)
%%% Author lane: Sanders + Gish only.
%%% Ported to public repo 2026-05-13 at
%%%   05_papers/number_theory/J03/manuscript/manuscript.tex
%%% ============================================================

\documentclass[11pt,reqno]{amsart}

\usepackage{amsmath,amssymb,amsthm,mathtools}
\usepackage[margin=1in]{geometry}
\usepackage[colorlinks=true,linkcolor=blue,citecolor=blue,urlcolor=blue]{hyperref}
\usepackage{microtype}
\usepackage{enumitem}

%% -------- theorem environments --------
\theoremstyle{plain}
\newtheorem{theorem}{Theorem}
\newtheorem{lemma}[theorem]{Lemma}
\newtheorem{proposition}[theorem]{Proposition}
\newtheorem{corollary}[theorem]{Corollary}

\theoremstyle{definition}
\newtheorem{definition}[theorem]{Definition}

\theoremstyle{remark}
\newtheorem*{remark}{Remark}

%% -------- macros --------
\DeclareMathOperator{\spf}{spf}
\DeclareMathOperator{\sinc}{sinc}
\newcommand{\Z}{\mathbb{Z}}
\newcommand{\N}{\mathbb{N}}
\newcommand{\R}{\mathbb{R}}
\newcommand{\Prm}{\mathbb{P}}
\newcommand{\Zb}{\Z/b\Z}

%% -------- title matter --------
\title[The First-G event and a discrete sinc${}^{2}$ identity]%
{The First-G Event and a Discrete Sinc${}^{2}$ Identity}

\author{B.~R.~Sanders}
\address{7Site LLC, Hot Springs, Arkansas, USA}
\email{brayden@7site.co}

\author{M.~Gish}
\address{Independent Researcher, Hot Springs, Arkansas, USA}
\email{monica.gish1992@gmail.com}

\date{\today}

\keywords{coprimality partition, smallest prime factor, Euler totient,
sieve of Eratosthenes, Fej\'{e}r kernel, sinc function, discrete Fourier
identity, rectangular pulse spectrum}

\subjclass[2020]{11A41, 11N05, 11A51, 11Y05, 42A16}

\begin{document}

\begin{abstract}
For an integer $b > 1$ and a growing alphabet $\{1,\dots,k\}$, the
\emph{coprimality partition} of the alphabet relative to $b$ separates
the units $C_{k}(b) = \{x : \gcd(x,b) = 1\}$ from the obstructions
$G_{k}(b) = \{1,\dots,k\}\setminus C_{k}(b)$. We establish a spectral
characterization of the obstruction sequence by means of a discrete
Fej\'{e}r-type quotient
$R(k,f) = \sin^{2}(\pi k/f)/(k^{2}\sin^{2}(\pi/f))$. The
\emph{spectral product} $f_{b}(k) := \prod_{p\mid b,\, p \text{ prime}}
R(k,p)$, defined as the product over the distinct prime factors of
$b$, vanishes at an integer $k \ge 1$ if and only if $\gcd(k,b) > 1$
(Theorem~\ref{thm:obstruction-zero}). The integer zero set of $f_{b}$
is therefore exactly the union of arithmetic progressions
$\bigcup_{p\mid b} p\,\N$, and $f_{b}$ acts as a continuous-in-$k$
indicator function for the obstruction event. The smallest zero
localizes at $k = \spf(b)$, recovering the First-G localization
theorem $k^{\star}(b) = \spf(b)$ as a corollary
(Theorem~\ref{thm:sync}); the asymptotic zero density of $f_{b}$
recovers the Euler product $1 - \varphi(\mathrm{rad}(b))/\mathrm{rad}(b)$
(Corollary~\ref{cor:asymp-density}). The closed form for $R(k,f)$
itself is the standard Fej\'{e}r-type calculation, included for
completeness; the synchronization at $k = \spf(b)$ and the spectral
characterization of the full obstruction sequence are the
contributions of this note. The continuum limit
$R(k,f) \to \sinc^{2}(k/f)$ as $f \to \infty$ is recorded as
context~(Theorem~\ref{thm:limit}).
The proofs are elementary throughout. Verification includes every
squarefree $b$ in $[2,500]$ (22{,}367 $(b,k)$ pairs, zero
counterexamples) and every prime $f$ in $[3,23]$ together with the
spectral-product zero locus over squarefree $b \in [2,100]$ at
exact integer arithmetic, with zero exceptions.
\end{abstract}

\maketitle

%% -------- section 0: lens and substrate (per SAVE_PLAN J03 §5) --------
\section*{Lens and substrate}\label{sec:lens}

This note works on $\Z$ (no specialized substrate) with the standard
objects: smallest prime factor $\spf(b)$, the discrete Fej\'{e}r
quotient $R(k, f)$, and the continuum $\sinc^{2}$ function. The objects
are not derived from a framework reading; they are the elementary
discrete-Fourier and divisibility structures of the integers. The
substantive contribution is the spectral characterization of the
obstruction sequence (Theorem~\ref{thm:obstruction-zero}); the
synchronization theorem (Theorem~\ref{thm:sync}) is the special case
at the smallest spectral zero.
We choose the harmonic-content presentation (smallest prime factor
$\times$ discrete Fej\'{e}r quotient on $\Z$) over a pure
$\Z/N\Z$-algebraic framing because the obstruction-zero correspondence
is a statement about integers and their prime divisibility, with the
spectral product $f_{b}$ supplying a single naturally-defined Fourier
indicator for $\gcd(k, b) > 1$. The natural setting is $\Z$, not
$\Z/N\Z$ with operator-labeled tables; no magma structure is required
for the present results, and the divisibility lattice of
$\mathrm{rad}(b)$ enters directly through the zero loci of the factors
$R(\cdot, p_{j})$. Subsequent companion papers in our program (the
Sanders--Gish $\sigma$-rate paper and the four-core fusion-closure
paper~\cite{SandersGish2026Sigma,SandersGish2026FourCore}) work over
$\Z/N\Z$ with operator-labeled tables; the present paper does not
require any of that machinery and stands alone.

%% -------- section 1: introduction --------
\section{Introduction}\label{sec:intro}

The sieve of Eratosthenes marks, for each integer $b > 1$ and each
alphabet size $k$, the elements of $\{1,\dots,k\}$ that share a prime
factor with $b$. Following standard usage \cite{Apostol1976,
HardyWright2008}, we organize the marking around the alphabet-size
parameter $k$ via the coprimality partition
\[
  C_{k}(b) = \{x\in\{1,\dots,k\}:\gcd(x,b)=1\},
  \qquad
  G_{k}(b) = \{1,\dots,k\}\setminus C_{k}(b),
\]
and call the smallest $k$ with $|G_{k}(b)|>0$ the \emph{First-G event}
of $b$. Section~\ref{sec:firstg} establishes that this $k$ is exactly
the smallest prime factor $p_{1} = \spf(b)$
(Theorem~\ref{thm:firstg}); the proof is one line of gcd reasoning.

Section~\ref{sec:harmonic} treats the harmonic counterpart. For an
integer parameter $f \ge 2$, we form the unit-alphabet exponential
sum at frequency $1/f$,
\[
  S(k,f) = \frac{1}{k}\sum_{j=1}^{k} e^{2\pi i j/f},
\]
and study its squared modulus $R(k,f) = |S(k,f)|^{2}$. This is a
discrete analogue of the Fej\'{e}r kernel \cite{Fejer1900,Zygmund2002},
the squared magnitude of a normalized finite exponential sum. We show
(Theorem~\ref{thm:closedform}) that $R(k,f)$ admits the closed form
\[
  R(k,f) = \frac{\sin^{2}(\pi k/f)}{k^{2}\sin^{2}(\pi/f)},
\]
and we record three immediate corollaries: $R(k,f)$ is strictly
decreasing on $\{1,\dots,f-1\}$, takes the value $1/(f-1)^{2}$ at
$k=f-1$, and vanishes at $k=f$.

Section~\ref{sec:sync} states the synchronization theorem
(Theorem~\ref{thm:sync}): for $f = p_{1} = \spf(b)$, the First-G
event of Theorem~\ref{thm:firstg} occurs at the same alphabet-size
value $k = p_{1}$ as the first integer zero of $R(k,p_{1})$ from
Theorem~\ref{thm:closedform}. The two events colocalize.

Section~\ref{sec:limit} treats the continuum limit: as $f\to\infty$
with $t = k/f$ fixed, $R(k,f) \to \sinc^{2}(t)$ where
$\sinc(t) = \sin(\pi t)/(\pi t)$ is the normalized sinc function
\cite{Shannon1949,OppenheimSchafer2010}. This identifies $R(k,f)$
as a discrete approximant to the rectangular-pulse power spectrum.

\paragraph{Tier discipline.} This paper PROVES (Theorem~\ref{thm:sync})
the synchronization of the First-G arithmetic event with the first
integer zero of $R(\cdot, \spf(b))$, and PROVES
(Theorem~\ref{thm:limit}) that $R(k,f) \to \sinc^{2}(k/f)$ as
$f \to \infty$. We COMPUTE the verifications: $22{,}367$ $(b,k)$ pairs
over squarefree $b \le 500$ for the First-G localization (script
\texttt{proof\_first\_g\_event.py}, runtime under $3$ s), and
machine-precision evaluation of the closed form for $8$ primes
(maximum deviation $4.44 \times 10^{-16}$, script
\texttt{verify\_J03.py}). The exact identity
$\sinc^{2}(1/2) = (2/3)/\zeta(2)$ is cited as STRUCTURAL RHYME, not
theorem --- it follows in one line from $\zeta(2) = \pi^{2}/6$. The
theorem-shaped question of why the corridor midpoint at $t = 1/2$
makes this identity structurally relevant is OPEN and is the natural
next paper.

\paragraph{Positioning and contribution.} The First-G localization
(Theorem~\ref{thm:firstg}) and the closed form for $R(k,f)$
(Theorem~\ref{thm:closedform}) are elementary in isolation: the first
unpacks the definition of $\spf(b)$ via a one-line gcd argument, and
the second is a Fej\'{e}r-type calculation in finite Fourier
analysis~\cite{Zygmund2002,OppenheimSchafer2010}. The
$\sinc^{2}$-continuum limit (Theorem~\ref{thm:limit}) is similarly a
direct asymptotic. The substantive content of the present note lies
in the \emph{spectral characterization of the obstruction sequence}
(Theorem~\ref{thm:obstruction-zero} and
Corollaries~\ref{cor:incl-excl},~\ref{cor:asymp-density}): for every
$b > 1$ with distinct prime factors $p_{1}, \dots, p_{r}$, the
product
\[
  f_{b}(k) = \prod_{j=1}^{r} R(k, p_{j})
\]
of discrete Fej\'{e}r quotients vanishes at $k \in \N$ precisely
when $\gcd(k, b) > 1$, and the synchronization theorem
(Theorem~\ref{thm:sync}) is the special case identifying the
\emph{smallest} such $k$ as $\spf(b)$. The spectral product $f_{b}$
acts as a continuous-in-$k$ indicator for the obstruction event, with
the divisibility lattice of $\mathrm{rad}(b)$ entering through the
zero loci of the factors. The asymptotic zero density of $f_{b}$
recovers the Euler product
$\varphi(\mathrm{rad}(b))/\mathrm{rad}(b)$
(Corollary~\ref{cor:asymp-density}). The contribution is therefore
the explicit Fourier-side encoding of the divisibility structure of
$b$ by a single, naturally-defined product spectral function, not
merely the synchronization at one value of $k$. Related classical
work on the distribution of the smallest prime factor
\cite{Erdos1959,Pomerance1985,Tenenbaum2015} and sieve-theoretic
treatments \cite{IwaniecKowalski2004,FriedlanderIwaniec2010} provide
the analytic-number-theory context. The present note also bridges
related work on operator-substrate composition tables over $\Zb$
\cite{SandersGish2026Sigma} and on dimensionless scalar-field models
with logarithmic potentials~\cite{SandersGishJohnson2026JCAP}.

\paragraph{Verification.} All numerical claims are verified by exact
computation. The First-G localization is checked for every squarefree
$b$ in $[2,500]$ ($305$ values, $22{,}367$ $(b,k)$ pairs, zero
counterexamples); Theorem~\ref{thm:firstg} itself does not require
squarefree-ness and the test set is chosen for convenience and
compatibility with the companion paper~\cite{SandersGish2026Sigma}.
The closed form for $R(k,f)$ is checked for every prime
$f \in \{3,5,7,11,13,17,19,23\}$ and every $k\in\{1,\dots,f+1\}$,
with maximum closed-form deviation below machine epsilon
($4.44\times 10^{-16}$). Verification scripts are documented in
Section~\ref{sec:verify}.

%% -------- section 2: setup --------
\section{Setup and Notation}\label{sec:setup}

Throughout, $b$ is an integer with $b > 1$. Write $b = p_{1}^{a_{1}}
\cdots p_{r}^{a_{r}}$ with $p_{1} < p_{2} < \cdots < p_{r}$ the
distinct prime factors of $b$, so that $p_{1} = \spf(b)$ is the
smallest prime factor.

\begin{definition}[Coprimality partition]\label{def:partition}
For $b > 1$ and $k \ge 1$, the coprimality partition of the alphabet
$\{1,\dots,k\}$ relative to $b$ is the ordered pair
\[
  C_{k}(b) = \{x \in \{1,\dots,k\} : \gcd(x,b) = 1\},
  \qquad
  G_{k}(b) = \{1,\dots,k\}\setminus C_{k}(b).
\]
The \emph{First-G event} of $b$ is the integer
\[
  k^{\star}(b) = \min\{k \ge 1 : |G_{k}(b)| > 0\}.
\]
\end{definition}

The map $b \mapsto k^{\star}(b)$ is well-defined for every $b > 1$
because $b \in G_{b}(b)$. The main theorem of
Section~\ref{sec:firstg} identifies $k^{\star}(b)$ explicitly.

\begin{definition}[Harmonic resonance]\label{def:resonance}
For an integer $f \ge 2$ and $k \ge 1$, define the unit-alphabet
exponential sum at frequency $1/f$ and its squared modulus by
\[
  S(k,f) = \frac{1}{k}\sum_{j=1}^{k} e^{2\pi i j/f},
  \qquad
  R(k,f) = |S(k,f)|^{2}.
\]
\end{definition}

The function $R(k,f)$ is the squared magnitude of the average of
the first $k$ terms of the geometric progression $e^{2\pi i j/f}$;
in finite Fourier analysis it is recognized as a discrete relative
of the Fej\'{e}r kernel~\cite{Fejer1900,Zygmund2002}, evaluated at
unit-alphabet sample sizes.

%% -------- section 3: First-G --------
\section{First-G Event Localization}\label{sec:firstg}

\begin{theorem}[First-G localization]\label{thm:firstg}
Let $b > 1$ be an integer with smallest prime factor $p_{1}=\spf(b)$.
Then $k^{\star}(b) = p_{1}$. More precisely:
\begin{enumerate}[label={\upshape(\roman*)}]
\item $|G_{k}(b)| = 0$ for every $k$ with $1 \le k < p_{1}$, and
\item $|G_{p_{1}}(b)| = 1$, with $G_{p_{1}}(b) = \{p_{1}\}$.
\end{enumerate}
\end{theorem}

\begin{proof}
\textit{(i)} Fix $k < p_{1}$ and $x \in \{1,\dots,k\}$. If
$\gcd(x,b) = d > 1$, then $d$ has a prime divisor $q$ with $q\mid b$;
hence $q \ge p_{1}$ by definition of $p_{1}$. But $q \mid x$ forces
$q \le x \le k < p_{1}$, a contradiction. Therefore $\gcd(x,b) = 1$
and $G_{k}(b) = \emptyset$.

\textit{(ii)} At $k = p_{1}$, the element $p_{1}$ enters the alphabet,
and $\gcd(p_{1},b) = p_{1} > 1$ since $p_{1}\mid b$, so
$p_{1}\in G_{p_{1}}(b)$. For any $x \in \{1,\dots,p_{1}\}$ with
$x \ne p_{1}$ we have $x < p_{1}$, so part (i) gives $\gcd(x,b)=1$.
Thus $G_{p_{1}}(b) = \{p_{1}\}$.
\end{proof}

\begin{remark}
The proof uses only the definitions of $p_{1}$ and $\gcd$.
Squarefree-ness of $b$ is not required. Although the proof goes
through verbatim for non-squarefree $b$ (e.g., for $b = 12$,
$\spf(b) = 2$ and $G_{2}(12) = \{2\}$ in the same way as
$G_{2}(6) = \{2\}$), we restrict the verification sample to squarefree
$b$ because that is the regime where the companion First-G application
program lives.
\end{remark}

\begin{corollary}[Stability window]\label{cor:stability}
For every $b > 1$ with $\spf(b) = p_{1}$, the interval
$\{1,\dots,p_{1}-1\}$ lies entirely in $C_{p_{1}-1}(b)$. The width of
this interval is $p_{1}-1$, and it is uniform across the family of
all integers with the same smallest prime factor.
\end{corollary}

\begin{corollary}[Phase-transition set]\label{cor:phase}
The image of $b\mapsto k^{\star}(b)$ over integers $b > 1$ is exactly
the set $\Prm$ of primes.
\end{corollary}

\begin{proof}
By Theorem~\ref{thm:firstg}, $k^{\star}(b) = \spf(b) \in \Prm$ for
every $b > 1$. Conversely, for any prime $p$, taking $b = p$ gives
$k^{\star}(p) = p$. Hence the image is $\Prm$.
\end{proof}

%% -------- section 4: closed form --------
\section{The Closed Form for $R(k,f)$}\label{sec:harmonic}

\begin{theorem}[Discrete sinc${}^{2}$ identity]\label{thm:closedform}
For every integer $f \ge 2$ and every integer $k \ge 1$,
\begin{equation}\label{eq:closedform}
  R(k,f) \;=\; \frac{\sin^{2}(\pi k/f)}{k^{2}\sin^{2}(\pi/f)}.
\end{equation}
\end{theorem}

\begin{proof}
The geometric-series formula gives
\[
  \sum_{j=1}^{k}e^{2\pi i j/f}
  \;=\; e^{2\pi i/f}\,\frac{1-e^{2\pi i k/f}}{1-e^{2\pi i/f}}.
\]
Taking the squared modulus of $(1/k)$ times this sum and applying
the elementary identity $|1-e^{i\theta}|^{2} = 4\sin^{2}(\theta/2)$,
\[
  R(k,f)
  = \frac{1}{k^{2}}\,
    \frac{|1-e^{2\pi i k/f}|^{2}}{|1-e^{2\pi i/f}|^{2}}
  = \frac{1}{k^{2}}\,
    \frac{4\sin^{2}(\pi k/f)}{4\sin^{2}(\pi/f)}
  = \frac{\sin^{2}(\pi k/f)}{k^{2}\sin^{2}(\pi/f)}.
\qedhere
\]
\end{proof}

\begin{corollary}[Endpoint values and monotonicity]\label{cor:endpoints}
For every prime $f \ge 2$:
\begin{enumerate}[label={\upshape(\roman*)}]
\item $R(1,f) = 1$ (full resonance at $k = 1$);
\item $R(f-1,f) = 1/(f-1)^{2}$, the last positive value of
      $R(\cdot,f)$ before the first zero at $k=f$;
\item $R(f,f) = 0$ (exact zero at $k=f$);
\item $R(k,f)$ is strictly decreasing on $\{1,\dots,f-1\}$.
\end{enumerate}
\end{corollary}

\begin{proof}
\textit{(i)}\ Direct from \eqref{eq:closedform}: at $k=1$,
$R(1,f) = \sin^{2}(\pi/f)/\sin^{2}(\pi/f) = 1$.

\textit{(iii)}\ At $k=f$: $\sin(\pi f/f) = \sin(\pi) = 0$, so
$R(f,f) = 0$ exactly.

\textit{(ii)}\ At $k=f-1$: $\sin(\pi(f-1)/f) = \sin(\pi - \pi/f)
= \sin(\pi/f)$, so $R(f-1,f) = \sin^{2}(\pi/f)/((f-1)^{2}
\sin^{2}(\pi/f)) = 1/(f-1)^{2}$.

\textit{(iv)}\ Factor the closed form \eqref{eq:closedform} as
\[
  R(k,f) = \left(\frac{\pi/f}{\sin(\pi/f)}\right)^{2}
           \left(\frac{\sin(\pi k/f)}{\pi k/f}\right)^{2}.
\]
The first factor is constant in $k$. For $1\le k\le f-1$ we have
$\pi k/f\in(0,\pi)$, and the function $x \mapsto \sin x/x$ is strictly
decreasing on $(0,\pi)$: its derivative is
$(x\cos x - \sin x)/x^{2}$, and the function
$h(x):= \sin x - x\cos x$ satisfies $h(0)=0$ and
$h'(x) = x\sin x > 0$ for $x\in(0,\pi)$, hence $h(x) > 0$ throughout
$(0,\pi)$, so $x\cos x - \sin x < 0$ and the derivative of
$\sin x/x$ is negative. The square of a positive
strictly decreasing function is strictly decreasing, so $R(k,f)$ is
strictly decreasing on $\{1,\dots,f-1\}$.
\end{proof}

\begin{remark}[Connection to Fej\'{e}r-type identities]\label{rem:fejer}
Writing $T_{k}(\theta) = \sum_{j=1}^{k}e^{ij\theta}$ for the
one-sided partial geometric sum evaluated at $\theta = 2\pi/f$,
\eqref{eq:closedform} is equivalent to
$|T_{k}(\theta)|^{2}/k^{2} = \sin^{2}(k\theta/2)/(k^{2}
\sin^{2}(\theta/2))$. This is the standard Fej\'{e}r-type
identity~\cite{Fejer1900,Zygmund2002}, specialized to rational
arguments $\theta/2\pi = 1/f$.
\end{remark}

%% -------- section 5: synchronization --------
\section{Synchronization of the First-G Event and the First Integer Zero of $R$}\label{sec:sync}

The First-G event of Theorem~\ref{thm:firstg} and the first integer
zero of $R(k,f)$ from Corollary~\ref{cor:endpoints} coincide for the
canonical choice $f = p_{1} = \spf(b)$.

\begin{theorem}[Synchronization]\label{thm:sync}
For every integer $b > 1$ with smallest prime factor $p_{1}=\spf(b)$,
the First-G event $k^{\star}(b)$ and the first integer zero of
$R(k,p_{1})$ in the range $\{1,\dots,p_{1}\}$ coincide:
\[
  k^{\star}(b) \;=\; \min\{k \in \{1,\dots,p_{1}\} : R(k,p_{1}) = 0\}
  \;=\; p_{1}.
\]
\end{theorem}

\begin{proof}
Theorem~\ref{thm:firstg} gives $k^{\star}(b) = p_{1}$. By
Theorem~\ref{thm:closedform},
\[
  R(k,p_{1}) = \frac{\sin^{2}(\pi k/p_{1})}{k^{2}\sin^{2}(\pi/p_{1})}.
\]
For $1 \le k < p_{1}$, $p_{1}$ is prime and $p_{1}\nmid k$, so
$\pi k/p_{1}\notin\pi\Z$ and $\sin(\pi k/p_{1})\neq 0$; hence
$R(k,p_{1}) > 0$. At $k = p_{1}$, $\sin(\pi) = 0$, so
$R(p_{1},p_{1}) = 0$. The first integer zero of $R(\cdot,p_{1})$ in
$\{1,\dots,p_{1}\}$ therefore occurs at $k = p_{1}$, equal to
$k^{\star}(b)$.
\end{proof}

The synchronization is the joint observation that a single integer
$k = p_{1}$ simultaneously witnesses two phenomena: the arithmetic
event ``the alphabet first contains an obstruction modulo $b$'' and
the harmonic event ``the unit-alphabet exponential sum at frequency
$1/p_{1}$ first vanishes.''

The synchronization at $k = p_{1}$ extends to the entire obstruction
sequence via a product spectral function, which we now establish.

\begin{definition}[Spectral product]\label{def:spectral-product}
Let $b > 1$ be an integer with distinct prime factors $p_{1} < p_{2}
< \cdots < p_{r}$. The \emph{spectral product} of $b$ is
\begin{equation}\label{eq:spectral-product}
  f_{b}(k) \;:=\; \prod_{j=1}^{r} R(k, p_{j}), \qquad k \ge 1.
\end{equation}
By Theorem~\ref{thm:closedform} this has the closed form
\[
  f_{b}(k) \;=\; \frac{1}{k^{2r}}\,
                 \prod_{j=1}^{r}\frac{\sin^{2}(\pi k/p_{j})}
                                    {\sin^{2}(\pi/p_{j})}.
\]
\end{definition}

\begin{theorem}[Obstruction--zero correspondence]\label{thm:obstruction-zero}
Let $b > 1$ with distinct prime factors $p_{1} < \cdots < p_{r}$, and
let $f_{b}$ be the spectral product of Definition~\ref{def:spectral-product}.
Then for every integer $k \ge 1$,
\begin{equation}\label{eq:obstruction-zero}
  f_{b}(k) \;=\; 0
  \quad\Longleftrightarrow\quad
  \gcd(k, b) > 1
  \quad\Longleftrightarrow\quad
  k \in G_{k}(b).
\end{equation}
Equivalently, the zero set of $f_{b}$ in $\N$ is
\[
  Z(f_{b}) \;=\; \bigcup_{j=1}^{r} p_{j}\,\N,
\]
the union of the arithmetic progressions of multiples of each
distinct prime factor of $b$.
\end{theorem}

\begin{proof}
By Theorem~\ref{thm:closedform},
$R(k, p_{j}) = \sin^{2}(\pi k/p_{j})/(k^{2}\sin^{2}(\pi/p_{j}))$.
For every $j$ the denominator $k^{2}\sin^{2}(\pi/p_{j})$ is strictly
positive (since $k \ge 1$ and $p_{j} \ge 2$), so
\[
  R(k, p_{j}) = 0
  \;\Longleftrightarrow\; \sin(\pi k/p_{j}) = 0
  \;\Longleftrightarrow\; p_{j} \mid k.
\]
The product $f_{b}(k)$ vanishes iff some factor $R(k, p_{j})$
vanishes iff $p_{j} \mid k$ for some $j$ iff $\gcd(k, b) > 1$
(since the distinct prime factors of $b$ are exactly the $p_{j}$).
The final equivalence $\gcd(k, b) > 1 \Leftrightarrow k \in G_{k}(b)$
is Definition~\ref{def:partition}.
\end{proof}

\begin{remark}\label{rem:5.1-as-special-case}
Theorem~\ref{thm:sync} is the special case of
Theorem~\ref{thm:obstruction-zero} at the smallest zero of $f_{b}$.
The smallest $k \ge 1$ with $f_{b}(k) = 0$ is the smallest element
of $\bigcup_{j} p_{j}\N$, which is $p_{1} = \spf(b)$. The
synchronization theorem therefore localizes the \emph{first}
spectral zero; Theorem~\ref{thm:obstruction-zero} characterizes
\emph{all} spectral zeros as obstruction events.
\end{remark}

\begin{remark}[No squarefree restriction; $f_{b}$ depends only on $\mathrm{rad}(b)$]\label{rem:radical}
Theorem~\ref{thm:obstruction-zero} holds for every integer $b > 1$,
not only for squarefree $b$. Definition~\ref{def:spectral-product}
forms the spectral product $f_{b}(k)$ from the \emph{distinct} prime
factors of $b$; these are precisely the prime factors of the radical
$\mathrm{rad}(b) = \prod_{p \mid b} p$, so $f_{b}(k) = f_{\mathrm{rad}(b)}(k)$
for every $k$. For example, $b = 12 = 2^{2}\cdot 3$ has distinct primes
$\{2, 3\}$, so $f_{12}(k) = R(k,2)R(k,3) = f_{6}(k)$, and the zero
set $\bigcup\{2\N, 3\N\}$ is the set of $k$ with $\gcd(k, 12) > 1$.
Equivalently, the equivalence $f_{b}(k) = 0 \Leftrightarrow \gcd(k, b) > 1$
factors through the radical: $\gcd(k, b) > 1$ iff $\gcd(k, \mathrm{rad}(b)) > 1$,
and the spectral side reads off the same prime divisibility condition.
Corollary~\ref{cor:asymp-density} is similarly already stated for
general $b > 1$, with the density $1 - \varphi(\mathrm{rad}(b))/\mathrm{rad}(b)$
depending only on the radical.
\end{remark}

\begin{corollary}[Inclusion--exclusion identity for $|G_{k}|$]\label{cor:incl-excl}
With $b$ and $f_{b}$ as in Theorem~\ref{thm:obstruction-zero}, the
number of obstructions in $\{1, \dots, k\}$ is
\[
  |G_{k}(b)|
  \;=\; \#\{\,j \in \{1, \dots, k\} : f_{b}(j) = 0\,\}
  \;=\; -\sum_{\substack{d \mid \mathrm{rad}(b)\\ d > 1}}
              \mu(d)\,\lfloor k/d \rfloor,
\]
where $\mathrm{rad}(b) = p_{1}\cdots p_{r}$ is the radical of $b$.
\end{corollary}

\begin{proof}
The first equality is Theorem~\ref{thm:obstruction-zero}. The second
is the standard inclusion-exclusion count of integers in $[1, k]$
divisible by at least one $p_{j}$, applied to the squarefree integer
$\mathrm{rad}(b)$ (whose divisors $d$ are squarefree products of the
$p_{j}$, and which has $\mu(d) = (-1)^{\omega(d)}$ where $\omega(d)$
is the number of distinct prime factors).
\end{proof}

\begin{corollary}[Asymptotic zero density of $f_{b}$]\label{cor:asymp-density}
For every $b > 1$ with distinct prime factors $p_{1}, \dots, p_{r}$,
\begin{equation}\label{eq:asymp-density}
  \lim_{k \to \infty}
  \frac{1}{k}\,\#\{\,j \in \{1, \dots, k\} : f_{b}(j) = 0\,\}
  \;=\; 1 - \prod_{j=1}^{r}\!\left(1 - \frac{1}{p_{j}}\right)
  \;=\; 1 - \frac{\varphi(\mathrm{rad}(b))}{\mathrm{rad}(b)}.
\end{equation}
\end{corollary}

\begin{proof}
By Theorem~\ref{thm:obstruction-zero} and
Corollary~\ref{cor:incl-excl}, the LHS equals
$\lim_{k\to\infty} |G_{k}(b)|/k$. The set $G_{k}(b)$ is the
complement in $\{1, \dots, k\}$ of the set of integers coprime to
$\mathrm{rad}(b)$; the latter has multiplicative density
$\varphi(\mathrm{rad}(b))/\mathrm{rad}(b) =
\prod_{j}(1-1/p_{j})$ by the Euler product
\cite{HardyWright2008,Apostol1976}. Hence the asymptotic obstruction
density is $1 - \prod_{j}(1-1/p_{j})$.
\end{proof}

\begin{remark}[The contribution]\label{rem:contribution}
The classical content of Corollary~\ref{cor:incl-excl} (the
inclusion-exclusion count of obstructions) and
Corollary~\ref{cor:asymp-density} (the Euler-product density)
is well-known. The contribution of
Theorem~\ref{thm:obstruction-zero} is the \emph{spectral
characterization}: the obstruction sequence is identified with the
zero set of an explicit product of discrete Fej\'{e}r quotients,
giving a Fourier-side encoding of the divisibility structure of
$b$ that goes beyond the synchronization at a single $k$. The
spectral product $f_{b}$ acts as a continuous (in $k$) indicator
function for $\gcd(k, b) > 1$, with the divisibility lattice of
$\mathrm{rad}(b)$ entering through the zero loci of the factors
$R(\cdot, p_{j})$.
\end{remark}

%% -------- section 6: continuum limit --------
\section{The Continuum Limit and the Rectangular Pulse}\label{sec:limit}

\begin{theorem}[$\sinc^{2}$ continuum limit]\label{thm:limit}
Let $\sinc(t) = \sin(\pi t)/(\pi t)$ for $t \ne 0$ and $\sinc(0)=1$.
Then for every $t \in \R$ with $t > 0$,
\[
  \lim_{f\to\infty,\ k/f \to t}
  R(k,f) \;=\; \sinc^{2}(t).
\]
\end{theorem}

\begin{proof}
From \eqref{eq:closedform}, $R(k,f) = \sin^{2}(\pi k/f)/(k^{2}
\sin^{2}(\pi/f))$. As $f \to \infty$ with $k/f \to t$ fixed,
$k \cdot \sin(\pi/f) = (\pi/f)\cdot k \cdot (1 + O(1/f^{2}))
\to \pi t$, so $k^{2}\sin^{2}(\pi/f) \to \pi^{2} t^{2}$. The
numerator satisfies $\sin^{2}(\pi k/f) \to \sin^{2}(\pi t)$.
Hence
\[
  R(k,f) \;\to\; \frac{\sin^{2}(\pi t)}{\pi^{2}t^{2}} \;=\; \sinc^{2}(t).
  \qedhere
\]
\end{proof}

\begin{remark}[Rectangular pulse spectrum]\label{rem:rectangular}
The function $\sinc^{2}(t)$ is the power spectral density of the
rectangular pulse $\mathbf{1}_{[0,1]}$, by direct evaluation of the
Fourier transform~\cite{Shannon1949,OppenheimSchafer2010}.
Theorem~\ref{thm:limit} therefore identifies $R(k,f)$ as a discrete
approximant to the rectangular-pulse spectrum sampled at rational
arguments $k/f$. The first integer zero of the discrete approximant
$R(\cdot,f)$, which by Corollary~\ref{cor:endpoints}(iii) sits at
$k = f$, is the discrete analogue of the first zero of $\sinc^{2}$
at $t = 1$.
\end{remark}

\begin{remark}[Specific exact values; structural rhyme]\label{rem:exactvalues}
Two evaluations of $\sinc^{2}$ that arise naturally as continuum
limits of $R(k,f)$ at rational arguments:
\begin{align*}
  \sinc^{2}(1/2) &= \frac{4}{\pi^{2}} = \frac{2/3}{\zeta(2)}, \\
  \sinc^{2}(1/10) &= \frac{25(\sqrt{5}-1)^{2}}{4\pi^{2}}.
\end{align*}
The first identity follows from $\sin(\pi/2) = 1$ together with
$\zeta(2) = \pi^{2}/6$ (one-line consequence). It is cited here as a
\emph{structural rhyme}, not as a theorem of the present paper: the
question of why the corridor midpoint at $t = 1/2$ makes this
identity structurally relevant is open. The second identity follows
from $\sin(\pi/10) = (\sqrt{5}-1)/4$. Discrete approximants
$R(k,f)$ with $k/f$ near these rational arguments converge to the
displayed exact values as $f \to \infty$.
\end{remark}

%% -------- section 7: verification --------
\section{Verification}\label{sec:verify}

The First-G localization (Theorem~\ref{thm:firstg}) is verified by
exact integer arithmetic for every squarefree $b$ in $[2,500]$ via
the script \texttt{proof\_first\_g\_event.py}. The total computational
load is $305$ squarefree $b$ values, $22{,}367$ $(b,k)$ pairs (each
$b$ contributes $\spf(b)$ pairs: $\spf(b)-1$ pairs for part (i) and
one pair for part (ii)), with zero counterexamples and runtime under
three seconds.

The closed form of Theorem~\ref{thm:closedform} is verified by
floating-point evaluation against the literal geometric sum for every
prime $f \in \{3,5,7,11,13,17,19,23\}$ and every $k \in
\{1,\dots,f+1\}$. The maximum deviation is $4.44 \times 10^{-16}$,
i.e., machine epsilon. The minimum-value identity $R(f-1,f) =
1/(f-1)^{2}$ from Corollary~\ref{cor:endpoints}(ii) is checked
directly:

\begin{center}
\begin{tabular}{c|c|c|c}
$f$ & $R(f-1,f)$ measured & $1/(f-1)^{2}$ & deviation\\
\hline
$3$ & $0.250000000$ & $0.250000000$ & $5.55\times 10^{-17}$\\
$5$ & $0.062500000$ & $0.062500000$ & $2.78\times 10^{-17}$\\
$7$ & $0.027777778$ & $0.027777778$ & $3.47\times 10^{-18}$\\
$11$ & $0.010000000$ & $0.010000000$ & $1.39\times 10^{-17}$\\
$13$ & $0.006944444$ & $0.006944444$ & $1.91\times 10^{-17}$\\
$17$ & $0.003906250$ & $0.003906250$ & $9.54\times 10^{-18}$\\
$19$ & $0.003086420$ & $0.003086420$ & $1.30\times 10^{-18}$\\
$23$ & $0.002066116$ & $0.002066116$ & $1.73\times 10^{-18}$\\
\end{tabular}
\end{center}

The synchronization theorem (Theorem~\ref{thm:sync}) is a direct
combination of the two preceding theorems, and the numerical checks
above are consistent with that synchronization.

%% -------- section 8: scope --------
\section{Scope and Limitations}\label{sec:scope}

We summarize what the present note does and does not claim.

\begin{enumerate}[label={\upshape(\arabic*)}]
\item The First-G event $k^{\star}(b)$ of Definition~\ref{def:partition}
      equals $\spf(b)$ for every $b > 1$
      (Theorem~\ref{thm:firstg}); this is an elementary one-line gcd
      argument and does not require squarefree-ness.

\item The closed form $R(k,f) = \sin^{2}(\pi k/f)/
      (k^{2}\sin^{2}(\pi/f))$ holds for every integer $f \ge 2$ and
      every integer $k \ge 1$ (Theorem~\ref{thm:closedform}). The
      proof is the standard Fej\'{e}r-type calculation in finite
      Fourier analysis~\cite{Fejer1900,Zygmund2002}; we make no
      claim of novelty for the identity itself, only for its
      synchronization with the First-G event.

\item The synchronization (Theorem~\ref{thm:sync}) is stated for
      $f = \spf(b)$. We do not claim a synchronization for other
      values of $f$, and we do not claim that the
      synchronization characterizes $b$ beyond its smallest prime
      factor.

\item The continuum limit (Theorem~\ref{thm:limit}) is a direct
      asymptotic. We do not claim any analytic-number-theory
      consequence: no statement about the distribution of primes in
      arithmetic progressions, no zero-free region for $\zeta(s)$,
      no consequence for the Riemann hypothesis, and no implication
      for cryptographic factoring difficulty.

\item The verification in Section~\ref{sec:verify} is exhaustive on
      the stated finite ranges. The theorems themselves are proved
      for all relevant $b$ and $f$, so the finite verification is
      a confirmation, not a pillar of the result.
\end{enumerate}

%% -------- references --------
\begin{thebibliography}{99}

\bibitem{Apostol1976}
T.~M.~Apostol.
\newblock \emph{Introduction to Analytic Number Theory}.
\newblock Undergraduate Texts in Mathematics. Springer, New York, 1976.

\bibitem{Erdos1959}
P.~Erd\H{o}s.
\newblock On the distribution of the smallest prime factor of $n$.
\newblock \emph{Mathematika}, 6:1--6, 1959.

\bibitem{Fejer1900}
L.~Fej\'{e}r.
\newblock Sur les fonctions born\'{e}es et int\'{e}grables.
\newblock \emph{Comptes Rendus de l'Acad\'{e}mie des Sciences, Paris},
131:984--987, 1900.

\bibitem{FriedlanderIwaniec2010}
J.~Friedlander and H.~Iwaniec.
\newblock \emph{Opera de Cribro}.
\newblock American Mathematical Society Colloquium Publications 57.
\newblock American Mathematical Society, Providence, RI, 2010.

\bibitem{HardyWright2008}
G.~H.~Hardy and E.~M.~Wright.
\newblock \emph{An Introduction to the Theory of Numbers}, 6th edition.
\newblock Revised by D.~R.~Heath-Brown and J.~H.~Silverman.
\newblock Oxford University Press, Oxford, 2008.

\bibitem{IwaniecKowalski2004}
H.~Iwaniec and E.~Kowalski.
\newblock \emph{Analytic Number Theory}.
\newblock American Mathematical Society Colloquium Publications 53.
\newblock American Mathematical Society, Providence, RI, 2004.

\bibitem{Lang2002}
S.~Lang.
\newblock \emph{Algebra}, 3rd edition.
\newblock Graduate Texts in Mathematics 211. Springer, New York, 2002.

\bibitem{OppenheimSchafer2010}
A.~V.~Oppenheim and R.~W.~Schafer.
\newblock \emph{Discrete-Time Signal Processing}, 3rd edition.
\newblock Prentice Hall, Upper Saddle River, NJ, 2010.

\bibitem{Pomerance1985}
C.~Pomerance.
\newblock On the distribution of round numbers.
\newblock In \emph{Number Theory (Ootacamund 1984)}, Lecture Notes in
Mathematics 1122, pages 173--200. Springer, Berlin, 1985.

\bibitem{Shannon1949}
C.~E.~Shannon.
\newblock Communication in the presence of noise.
\newblock \emph{Proceedings of the IRE}, 37(1):10--21, 1949.

\bibitem{Tenenbaum2015}
G.~Tenenbaum.
\newblock \emph{Introduction to Analytic and Probabilistic Number Theory},
3rd edition.
\newblock Cambridge Studies in Advanced Mathematics 46. Cambridge
University Press, Cambridge, 2015.

\bibitem{Zygmund2002}
A.~Zygmund.
\newblock \emph{Trigonometric Series}, 3rd edition.
\newblock With a foreword by R.~A.~Fefferman.
Cambridge University Press, Cambridge, 2002.

%% -- companion papers ------------------------
\bibitem{SandersGish2026Sigma}
B.~R.~Sanders and M.~Gish.
\newblock Non-Associativity Decay in Binary Composition Tables over
$\mathbb{Z}/N\mathbb{Z}$.
\newblock Preprint, 2026.

\bibitem{SandersGish2026FourCore}
B.~R.~Sanders and M.~Gish.
\newblock Joint Closure and a Closed-Form Algebraic Attractor of Two
Commutative Binary Operations on $\mathbb{Z}/10\mathbb{Z}$.
\newblock Preprint, 2026.

\bibitem{SandersGishJohnson2026JCAP}
B.~R.~Sanders, M.~Gish, and H.~J.~Johnson.
\newblock Logarithmic Quintessence: A Dimensionless Scalar Dark Energy
Model with an Analytic Vacuum.
\newblock Preprint, 2026.

\end{thebibliography}

\end{document}
